% Computer Science A --- Angular (practice bank).
% Formatted from raw notes; TypeScript items are in \texttt{web-typescript.tex}.

\runningheader{Computer Science A}{Web: Angular}{Page \thepage\ of \numpages}

\begingradingrange{csawebadvng}

\uplevel{\par\textbf{Angular framework}\par}

\question[4]
Which lifecycle hook is designed for cleanup tasks such as unsubscribing from Observables? \emph{Select all that apply.}
\begin{checkboxes}
  \CorrectChoice \lstinline[style=tsinline]|ngOnDestroy|
  \choice \lstinline[style=tsinline]|ngOnInit|
  \choice \lstinline[style=tsinline]|ngOnChanges|
  \choice \lstinline[style=tsinline]|ngAfterViewInit|
\end{checkboxes}
\begin{solution}
  \lstinline[style=tsinline]|ngOnDestroy| runs immediately before Angular destroys a component or directive. Unsubscribe from long-lived Observables, clear timers, and detach manual listeners here to avoid memory leaks. \lstinline[style=tsinline]|ngOnInit| initializes after inputs are set; \lstinline[style=tsinline]|ngOnChanges| reacts to \lstinline[style=tsinline]|@Input()| updates; \lstinline[style=tsinline]|ngAfterViewInit| runs after the view (and child views) exist.
\begin{lstlisting}[style=ts]
export class UserListComponent implements OnInit, OnDestroy {
  private sub?: Subscription;

  ngOnInit() {
    this.sub = this.users$.subscribe();
  }

  ngOnDestroy() {
    this.sub?.unsubscribe();
  }
}
\end{lstlisting}
\end{solution}

\question[4]
TRUE or FALSE: In Angular templates, property binding uses the syntax \lstinline[style=htmlinline]|(property)="value"|.
\begin{choices}
  \choice TRUE
  \CorrectChoice FALSE
\end{choices}
\begin{solution}
  \textbf{Property binding} flows data from the component class into the template with square brackets: \lstinline[style=htmlinline]|[property]="value"|. Parentheses \lstinline[style=htmlinline]|(event)="handler()"| are \textbf{event binding} (template to class). Two-way binding combines both: \lstinline[style=htmlinline]|[(ngModel)]="value"|.
\begin{lstlisting}[style=htmlinline]
<img [src]="photoUrl" alt="Profile">
<button (click)="save()">Save</button>
<input [(ngModel)]="title">
\end{lstlisting}
\end{solution}

\question[4]
TRUE or FALSE: The Angular CLI command \lstinline[style=bashinline]|ng serve| compiles code in real time and can live-reload the browser.
\begin{choices}
  \CorrectChoice TRUE
  \choice FALSE
\end{choices}
\begin{solution}
  \lstinline[style=bashinline]|ng serve| starts a development server, watches project files, recompiles on save, and refreshes or hot-updates the app so changes appear without a manual rebuild step.
\begin{lstlisting}[style=bash]
ng serve          # dev server, default http://localhost:4200
ng serve --open   # compile, serve, and open the browser
\end{lstlisting}
\end{solution}

\question[4]
To pass data from a parent component down to a child component, which decorator is used on the child? \emph{Select all that apply.}
\begin{checkboxes}
  \choice \lstinline[style=tsinline]|@Component()|
  \choice \lstinline[style=tsinline]|@Injectable()|
  \choice \lstinline[style=tsinline]|@Output()|
  \CorrectChoice \lstinline[style=tsinline]|@Input()|
\end{checkboxes}
\begin{solution}
  \lstinline[style=tsinline]|@Input()| marks a child property that the parent sets via property binding. \lstinline[style=tsinline]|@Output()| emits events upward with an \lstinline[style=tsinline]|EventEmitter|. \lstinline[style=tsinline]|@Component()| declares the component itself; \lstinline[style=tsinline]|@Injectable()| marks a service for dependency injection.
\begin{lstlisting}[style=ts]
@Component({ selector: 'app-child', template: `<p>{{ label }}</p>` })
export class ChildComponent {
  @Input() label = '';
}

// Parent template: <app-child [label]="parentTitle"></app-child>
\end{lstlisting}
\end{solution}

\question[4]
TRUE or FALSE: In an Angular single-page application (SPA), the browser requests a brand-new HTML document from the server on every route change.
\begin{choices}
  \choice TRUE
  \CorrectChoice FALSE
\end{choices}
\begin{solution}
  An SPA loads \texttt{index.html} once. Angular's client-side \textbf{Router} swaps components inside the existing page and updates the URL without a full server round-trip for a new HTML document, which keeps navigation fast.
\begin{lstlisting}[style=htmlinline]
<!-- index.html loaded once; <app-root> hosts routed views -->
<app-root></app-root>
\end{lstlisting}
\begin{lstlisting}[style=ts]
// Router config maps paths to components, not separate .html pages
{ path: 'labs', component: LabsComponent }
\end{lstlisting}
\end{solution}

\question[4]
TRUE or FALSE: The modern \lstinline[style=htmlinline]|@for| control flow requires a \lstinline[style=htmlinline]|track| expression for optimal rendering.
\begin{choices}
  \CorrectChoice TRUE
  \choice FALSE
\end{choices}
\begin{solution}
  In Angular's built-in control flow, \lstinline[style=htmlinline]|track| is \textbf{required}---omitting it is a compile error. The track expression (often a stable id) lets change detection identify added, removed, or reordered items instead of re-rendering the entire list, replacing the older \lstinline[style=htmlinline]|*ngFor| \lstinline[style=tsinline]|trackBy| pattern.
\begin{lstlisting}[style=htmlinline]
@for (item of items; track item.id) {
  <li>{{ item.name }}</li>
} @empty {
  <li>No items yet.</li>
}
\end{lstlisting}
\end{solution}

\question[4]
What is the correct syntax for two-way data binding in an Angular template? \emph{Select all that apply.}
\begin{checkboxes}
  \choice \lstinline[style=htmlinline]|{{ value }}|
  \choice \lstinline[style=htmlinline]|[value]="data"|
  \choice \lstinline[style=htmlinline]|(value)="handler()"|
  \CorrectChoice \lstinline[style=htmlinline]|[(ngModel)]="data"|
\end{checkboxes}
\begin{solution}
  \lstinline[style=htmlinline]|[(ngModel)]| combines property binding \lstinline[style=htmlinline]|[ ]| (class to template) and event binding \lstinline[style=htmlinline]|( )| (template to class), keeping an input and a component property in sync. \lstinline[style=htmlinline]|{{ }}| is one-way interpolation; \lstinline[style=htmlinline]|[value]| is one-way into the DOM; \lstinline[style=htmlinline]|(value)| listens for events only.
\begin{lstlisting}[style=htmlinline]
<label>
  Title
  <input [(ngModel)]="title" name="title">
</label>
<p>Preview: {{ title }}</p>
\end{lstlisting}
\end{solution}

\endgradingrange{csawebadvng}
